\subchapter{Key management}{
	Objective: Experiment with PKCS\#11 and OP-TEE}

\section{Background}

The OP-TEE codebase includes a
\href{https://github.com/OP-TEE/optee_os/tree/master/ta/pkcs11}{PKCS\#11 Trusted Application}.

First, confirm that this TA is present on your system.

\section{Setting up the TA}

We could write a small client application as we did in the \textit{secure world} lab, but
we don't need to.
Indeed, the API exposed by the TA can be compared to the internal API of an HSM, but the main
advantage of PKCS\#11 is the high-level API specification, which lets manufacturers provide
compatibility modules in the form of shared objects.

OP-TEE is no exception, and the compatibility module is implemented by the
\href{https://optee.readthedocs.io/en/latest/building/gits/optee_client.html}{OP-TEE client},
in the form of \href{https://github.com/OP-TEE/optee_client/tree/master/libckteec}{libckteec},
as documented \href{https://optee.readthedocs.io/en/latest/building/userland_integration.html#pkcs-11-driver}{here}.

You can confirm that everything is working by showing some general info:
\begin{bashinput}
root@freiheit93:~# p11 -I
Cryptoki version 2.40
Manufacturer     Linaro
Library          OP-TEE PKCS11 Cryptoki library (ver 0.1)
Using slot 0 with a present token (0x0)
\end{bashinput}
Notice the "Cryptoki" name showing up.

The next steps are the same as for a "normal" HSM, and described in the documenation linked above:
\begin{itemize}
\item Initialize the token, with a label and a Security Office (SO) PIN:
\begin{bashinput}
root@freiheit93:~# p11 --init-token --label optee
Using slot 0 with a present token (0x0)
Please enter the new SO PIN: 
Please enter the new SO PIN (again): 
Token successfully initialized

root@freiheit93:~# p11 -L
Available slots:
Slot 0 (0x0): OP-TEE PKCS11 TA - TEE UUID 94e9ab89-4c43-56ea-8b35-45dc07226830
  token label        : optee
  token manufacturer : Linaro
  token model        : OP-TEE TA
  token flags        : login required, rng, token initialized
  hardware version   : 0.0
  firmware version   : 0.1
  serial num         : 0000000000000000
  pin min/max        : 4/128
Slot 1 (0x1): OP-TEE PKCS11 TA - TEE UUID 94e9ab89-4c43-56ea-8b35-45dc07226830
  token state:   uninitialized
Slot 2 (0x2): OP-TEE PKCS11 TA - TEE UUID 94e9ab89-4c43-56ea-8b35-45dc07226830
  token state:   uninitialized

\end{bashinput}
\item Log in as the Security Office and create a User PIN:
\begin{bashinput}
# p11 -l --login-type so --init-pin
Using slot 0 with a present token (0x0)
Logging in to "optee".
Please enter SO PIN: 
Please enter the new PIN: 
Please enter the new PIN again: 
User PIN successfully initialized
\end{bashinput}
\end{itemize}


\section{Using the TA as a soft HSM}

Now that the (virtual) token has been initialized, we can use it to store the target's private key
that we generated earlier. This is done using the \code{--write-object} command:
\begin{bashinput}
# p11 -l --login-type user --write-object target.key --type privkey --label 'Target Key'
Using slot 0 with a present token (0x0)
Logging in to "optee".
Please enter User PIN: 
Created private key:
Private Key Object; RSA 
  label:      Target Key
  Usage:      none
  Access:     sensitive
\end{bashinput}

Unfortunately, this key risks to not be very useful. Experiment with the commands to fix that,
and please make sure that key cannot be used to sign MD5 hashes.

For the purpose of the first lab, we generate the target key using \code{openssl}. This is typically
the type of task that an HSM is well suited for.

Generate a new RSA 4096 key directly in OP-TEE.

\section{Thinking about HW storage}

Once your key is properly setup, reboot the device, and list the keys in the slot you were using:

\begin{bashinput}
$ p11 -l -O
\end{bashinput}

You should still see the keys that you created in the previous step. But OP-TEE itself does not
have a filesystem. Looking at the OP-TEE documentation and sources, try and figure out how
OP-TEE stored those keys.

We now have an RSA key that is tied to the hardware, and cannot leave the secure world in cleartext.
However, this key is not yet integrated into our PKI. We're going to enroll it now, following
similar steps as we did earlier.

In order to be able to use this key with OpenSSL, however, we need specific configuration:

\begin{bashinput}
$ cat /root/openssl.conf 
openssl_conf = openssl_init

[openssl_init]
engines = engine_section

[engine_section]
pkcs11 = pkcs11_section

[pkcs11_section]
engine_id = pkcs11
dynamic_path = /usr/lib/engines-3/pkcs11.so
MODULE_PATH = /usr/lib/libckteec.so.0.1.0
PIN = 5678

\end{bashinput}

Now we can generate a CSR:

\begin{bashinput}
# OPENSSL_CONF=openssl.conf openssl req -new \
	-engine pkcs11 \
	-keyform engine \
	-key pkcs11:object=TargetOnlyKey \
	-out target.csr
\end{bashinput}

After this, we can follow the steps from the first lab to sign the CSR.
